% --- Archon physics formalization source begin ---
% archon:physics
% archon:covers PhyXMiniProblems/problem_phyx_mini_0471.lean
% archon:source-report reports/phyx_mini/problem_phyx_mini_0471.source.json

\chapter{Physics problem phyx\_mini\_0471}
\label{ch:PhyXMiniProblems_problem_phyx_mini_0471}

\paragraph{Problem source.}
\#\# Physical scenario

The pV-diagram of figure shows a series of thermodynamic processes. In process ab, 150 J of heat is added to the system; in process bd, 600 J of heat is added.

\#\# Question

Find the total heat added in process acd.

\#\# Answer choices

- A: A: 510J

- B: B: 600J

- C: C: 490J

- D: D: 550J

\#\# Figure caption (auxiliary; use the image as primary evidence)

The image is a pressure-volume (p-V) diagram depicting a cyclic process involving a gas. It features the following elements:

- **Axes**: 
  - The vertical axis represents pressure (p) in pascals (Pa).
  - The horizontal axis represents volume (V) in cubic meters (m³).

- **Points**:
  - Four points labeled \textbackslash{}(a\textbackslash{}), \textbackslash{}(b\textbackslash{}), \textbackslash{}(c\textbackslash{}), and \textbackslash{}(d\textbackslash{}) form a rectangle, suggesting a sequence of processes.
  - The coordinates of these points correspond to specific values on the p-V axes.

- **Pressure Values**:
  - At point \textbackslash{}(b\textbackslash{}) and \textbackslash{}(d\textbackslash{}), pressure is \textbackslash{}(8.0 \textbackslash{}times 10\textasciicircum{}4 \textbackslash{}, \textbackslash{}text\{Pa\}\textbackslash{}).
  - At point \textbackslash{}(a\textbackslash{}) and \textbackslash{}(c\textbackslash{}), pressure is \textbackslash{}(3.0 \textbackslash{}times 10\textasciicircum{}4 \textbackslash{}, \textbackslash{}text\{Pa\}\textbackslash{}).

- **Volume Values**:
  - At point \textbackslash{}(a\textbackslash{}) and \textbackslash{}(b\textbackslash{}), volume is \textbackslash{}(2.0 \textbackslash{}times 10\textasciicircum{}\{-3\} \textbackslash{}, \textbackslash{}text\{m\}\textasciicircum{}3\textbackslash{}).
  - At point \textbackslash{}(c\textbackslash{}) and \textbackslash{}(d\textbackslash{}), volume is \textbackslash{}(5.0 \textbackslash{}times 10\textasciicircum{}\{-3\} \textbackslash{}, \textbackslash{}text\{m\}\textasciicircum{}3\textbackslash{}).

- **Process Path**:
  - The path follows: \textbackslash{}(a \textbackslash{}rightarrow b\textbackslash{}) (horizontal right), \textbackslash{}(b \textbackslash{}rightarrow d\textbackslash{}) (vertical up), \textbackslash{}(d \textbackslash{}rightarrow c\textbackslash{}) (horizontal left), and \textbackslash{}(c \textbackslash{}rightarrow a\textbackslash{}) (vertical down).
  - Arrows indicate the direction of each process in the cycle. 

- **Origin**:
  - Point \textbackslash{}(O\textbackslash{}) is marked at the origin where pressure and volume are zero. 

The diagram illustrates the cyclic thermodynamic process, typically useful in understanding concepts like work done, heat transfer, and transformations in engines or other systems involving gases.

\#\# Dataset metadata

Laws of Thermodynamics; Multi-Formula Reasoning, Physical Model Grounding Reasoning

\paragraph{Recorded answer/context.}
B: B: 600J

\paragraph{Figure/image path.}
/root/proposal\_for\_physic/science-mango/phyx\_mini\_run/phyx\_data/test\_image/471.png

\paragraph{Formalization target.}
create a compiling Lean file with sorry bodies at `PhyXMiniProblems/problem\_phyx\_mini\_0471.lean`. The Lean declarations must preserve the physical quantities, dimensions or dimensional roles, figure labels, governing-law hypotheses, and final relation expressed by this problem.
Use Mathlib/Physlib names found through LeanExplore where available. If a physics API is missing, introduce faithful local abstractions rather than scalar placeholder aliases.

\begin{theorem}[Physics formalization target]
\label{thm:physics:phyx_mini_0471:target}
\lean{PhyXMiniProblems.ProblemPhyXMini0471.problem_phyx_mini_0471}
\uses{def:physics:phyx-mini-0471:phyxminiproblems-problemphyxmini0471-energyinjoules, def:physics:phyx-mini-0471:phyxminiproblems-problemphyxmini0471-rectangularpvprocesssetup, def:physics:phyx-mini-0471:phyxminiproblems-problemphyxmini0471-matchesclosedsystemscenario, def:physics:phyx-mini-0471:phyxminiproblems-problemphyxmini0471-matchesprimarypressurevolumefigure, def:physics:phyx-mini-0471:phyxminiproblems-problemphyxmini0471-matchesgivenheatreadouts, def:physics:phyx-mini-0471:phyxminiproblems-problemphyxmini0471-hasphysicalpressurevolumecoordinates, def:physics:phyx-mini-0471:phyxminiproblems-problemphyxmini0471-satisfiesquasistaticboundaryworklaw, def:physics:phyx-mini-0471:phyxminiproblems-problemphyxmini0471-satisfiespathworkadditivity, def:physics:phyx-mini-0471:phyxminiproblems-problemphyxmini0471-satisfiesclosedsystemfirstlaw, def:physics:phyx-mini-0471:phyxminiproblems-problemphyxmini0471-matchesdisplayedheat}
The heat added along path \(a\to c\to d\) is \(600\,\mathrm J\), answer B.
\end{theorem}
\begin{proof}
Leg \(ab\) is isochoric, so its zero work and \(150\,\mathrm J\) heat give
\(\Delta U_{ab}=150\,\mathrm J\).  On \(bd\),
\[
 W_{bd}=80000(0.005-0.002)=240\,\mathrm J,
\]
so its \(600\,\mathrm J\) heat gives \(\Delta U_{bd}=360\,\mathrm J\).
Thus \(U_d-U_a=510\,\mathrm J\).  Along the alternative path only \(ac\)
does work:
\[
 W_{acd}=30000(0.005-0.002)=90\,\mathrm J.
\]
The first law now gives \(Q_{acd}=510+90=600\,\mathrm J\), exactly the value
displayed by B.
\end{proof}

% --- Archon declaration topology begin ---
\section{Declaration topology}
\begin{definition}[Volume Quantity]
  \label{def:physics:phyx-mini-0471:phyxminiproblems-problemphyxmini0471-volumequantity}
  \lean{PhyXMiniProblems.ProblemPhyXMini0471.VolumeQuantity}
  A signed physical volume carrying dimension L³
\end{definition}
\begin{proof}
  This is the declared model definition of Volume Quantity.
\end{proof}
\begin{definition}[Pressure Quantity]
  \label{def:physics:phyx-mini-0471:phyxminiproblems-problemphyxmini0471-pressurequantity}
  \lean{PhyXMiniProblems.ProblemPhyXMini0471.PressureQuantity}
  A physical pressure, using Physlib's pressure dimension
\end{definition}
\begin{proof}
  This is the declared model definition of Pressure Quantity.
\end{proof}
\begin{definition}[Energy Quantity]
  \label{def:physics:phyx-mini-0471:phyxminiproblems-problemphyxmini0471-energyquantity}
  \lean{PhyXMiniProblems.ProblemPhyXMini0471.EnergyQuantity}
  Signed heat, work, or internal energy, using Physlib's energy dimension
\end{definition}
\begin{proof}
  This is the declared model definition of Energy Quantity.
\end{proof}
\begin{definition}[volume In Cubic Metres]
  \label{def:physics:phyx-mini-0471:phyxminiproblems-problemphyxmini0471-volumeincubicmetres}
  \lean{PhyXMiniProblems.ProblemPhyXMini0471.volumeInCubicMetres}
  \uses{def:physics:phyx-mini-0471:phyxminiproblems-problemphyxmini0471-volumequantity}
  Read a physical volume in coherent-SI cubic metres
\end{definition}
\begin{proof}
  This is the declared model definition of volume In Cubic Metres.
\end{proof}
\begin{definition}[pressure In Pascals]
  \label{def:physics:phyx-mini-0471:phyxminiproblems-problemphyxmini0471-pressureinpascals}
  \lean{PhyXMiniProblems.ProblemPhyXMini0471.pressureInPascals}
  \uses{def:physics:phyx-mini-0471:phyxminiproblems-problemphyxmini0471-pressurequantity}
  Read a physical pressure in coherent-SI pascals
\end{definition}
\begin{proof}
  This is the declared model definition of pressure In Pascals.
\end{proof}
\begin{definition}[energy In Joules]
  \label{def:physics:phyx-mini-0471:phyxminiproblems-problemphyxmini0471-energyinjoules}
  \lean{PhyXMiniProblems.ProblemPhyXMini0471.energyInJoules}
  \uses{def:physics:phyx-mini-0471:phyxminiproblems-problemphyxmini0471-energyquantity}
  Read physical heat, work, or internal energy in coherent-SI joules
\end{definition}
\begin{proof}
  This is the declared model definition of energy In Joules.
\end{proof}
\begin{definition}[State Label]
  \label{def:physics:phyx-mini-0471:phyxminiproblems-problemphyxmini0471-statelabel}
  \lean{PhyXMiniProblems.ProblemPhyXMini0471.StateLabel}
  The four state labels printed beside the plotted points
\end{definition}
\begin{proof}
  This is the declared model definition of State Label.
\end{proof}
\begin{definition}[Process Leg]
  \label{def:physics:phyx-mini-0471:phyxminiproblems-problemphyxmini0471-processleg}
  \lean{PhyXMiniProblems.ProblemPhyXMini0471.ProcessLeg}
  The four directed straight legs shown by blue arrows in the image
\end{definition}
\begin{proof}
  This is the declared model definition of Process Leg.
\end{proof}
\begin{definition}[initial State]
  \label{def:physics:phyx-mini-0471:phyxminiproblems-problemphyxmini0471-processleg-initialstate}
  \lean{PhyXMiniProblems.ProblemPhyXMini0471.ProcessLeg.initialState}
  \uses{def:physics:phyx-mini-0471:phyxminiproblems-problemphyxmini0471-statelabel, def:physics:phyx-mini-0471:phyxminiproblems-problemphyxmini0471-processleg}
  Initial state of each displayed directed leg
\end{definition}
\begin{proof}
  This is the declared model definition of initial State.
\end{proof}
\begin{definition}[final State]
  \label{def:physics:phyx-mini-0471:phyxminiproblems-problemphyxmini0471-processleg-finalstate}
  \lean{PhyXMiniProblems.ProblemPhyXMini0471.ProcessLeg.finalState}
  \uses{def:physics:phyx-mini-0471:phyxminiproblems-problemphyxmini0471-statelabel, def:physics:phyx-mini-0471:phyxminiproblems-problemphyxmini0471-processleg}
  Final state of each displayed directed leg
\end{definition}
\begin{proof}
  This is the declared model definition of final State.
\end{proof}
\begin{definition}[Process Path]
  \label{def:physics:phyx-mini-0471:phyxminiproblems-problemphyxmini0471-processpath}
  \lean{PhyXMiniProblems.ProblemPhyXMini0471.ProcessPath}
  The two complete two-leg routes from state a to state d
\end{definition}
\begin{proof}
  This is the declared model definition of Process Path.
\end{proof}
\begin{definition}[first Leg]
  \label{def:physics:phyx-mini-0471:phyxminiproblems-problemphyxmini0471-processpath-firstleg}
  \lean{PhyXMiniProblems.ProblemPhyXMini0471.ProcessPath.firstLeg}
  \uses{def:physics:phyx-mini-0471:phyxminiproblems-problemphyxmini0471-processleg, def:physics:phyx-mini-0471:phyxminiproblems-problemphyxmini0471-processpath}
  First directed leg of a complete route
\end{definition}
\begin{proof}
  This is the declared model definition of first Leg.
\end{proof}
\begin{definition}[second Leg]
  \label{def:physics:phyx-mini-0471:phyxminiproblems-problemphyxmini0471-processpath-secondleg}
  \lean{PhyXMiniProblems.ProblemPhyXMini0471.ProcessPath.secondLeg}
  \uses{def:physics:phyx-mini-0471:phyxminiproblems-problemphyxmini0471-processleg, def:physics:phyx-mini-0471:phyxminiproblems-problemphyxmini0471-processpath}
  Second directed leg of a complete route
\end{definition}
\begin{proof}
  This is the declared model definition of second Leg.
\end{proof}
\begin{definition}[initial State]
  \label{def:physics:phyx-mini-0471:phyxminiproblems-problemphyxmini0471-processpath-initialstate}
  \lean{PhyXMiniProblems.ProblemPhyXMini0471.ProcessPath.initialState}
  \uses{def:physics:phyx-mini-0471:phyxminiproblems-problemphyxmini0471-statelabel, def:physics:phyx-mini-0471:phyxminiproblems-problemphyxmini0471-processpath}
  Both complete routes begin at state a
\end{definition}
\begin{proof}
  This is the declared model definition of initial State.
\end{proof}
\begin{definition}[final State]
  \label{def:physics:phyx-mini-0471:phyxminiproblems-problemphyxmini0471-processpath-finalstate}
  \lean{PhyXMiniProblems.ProblemPhyXMini0471.ProcessPath.finalState}
  \uses{def:physics:phyx-mini-0471:phyxminiproblems-problemphyxmini0471-statelabel, def:physics:phyx-mini-0471:phyxminiproblems-problemphyxmini0471-processpath}
  Both complete routes end at state d
\end{definition}
\begin{proof}
  This is the declared model definition of final State.
\end{proof}
\begin{definition}[Process Kind]
  \label{def:physics:phyx-mini-0471:phyxminiproblems-problemphyxmini0471-processkind}
  \lean{PhyXMiniProblems.ProblemPhyXMini0471.ProcessKind}
  Thermodynamic constraint represented by a straight p--V segment
\end{definition}
\begin{proof}
  This is the declared model definition of Process Kind.
\end{proof}
\begin{definition}[Segment Orientation]
  \label{def:physics:phyx-mini-0471:phyxminiproblems-problemphyxmini0471-segmentorientation}
  \lean{PhyXMiniProblems.ProblemPhyXMini0471.SegmentOrientation}
  Orientation of a displayed segment in the pressure--volume plane
\end{definition}
\begin{proof}
  This is the declared model definition of Segment Orientation.
\end{proof}
\begin{definition}[Figure Axis]
  \label{def:physics:phyx-mini-0471:phyxminiproblems-problemphyxmini0471-figureaxis}
  \lean{PhyXMiniProblems.ProblemPhyXMini0471.FigureAxis}
  The two axes of the supplied plot
\end{definition}
\begin{proof}
  This is the declared model definition of Figure Axis.
\end{proof}
\begin{definition}[Axis Quantity]
  \label{def:physics:phyx-mini-0471:phyxminiproblems-problemphyxmini0471-axisquantity}
  \lean{PhyXMiniProblems.ProblemPhyXMini0471.AxisQuantity}
  Physical quantity assigned to an axis of the supplied plot
\end{definition}
\begin{proof}
  This is the declared model definition of Axis Quantity.
\end{proof}
\begin{definition}[Axis Display Unit]
  \label{def:physics:phyx-mini-0471:phyxminiproblems-problemphyxmini0471-axisdisplayunit}
  \lean{PhyXMiniProblems.ProblemPhyXMini0471.AxisDisplayUnit}
  Unit text carried by an axis of the supplied plot
\end{definition}
\begin{proof}
  This is the declared model definition of Axis Display Unit.
\end{proof}
\begin{definition}[Axis Symbol]
  \label{def:physics:phyx-mini-0471:phyxminiproblems-problemphyxmini0471-axissymbol}
  \lean{PhyXMiniProblems.ProblemPhyXMini0471.AxisSymbol}
  Mathematical symbol printed at an axis tip
\end{definition}
\begin{proof}
  This is the declared model definition of Axis Symbol.
\end{proof}
\begin{definition}[Pressure Volume Figure]
  \label{def:physics:phyx-mini-0471:phyxminiproblems-problemphyxmini0471-pressurevolumefigure}
  \lean{PhyXMiniProblems.ProblemPhyXMini0471.PressureVolumeFigure}
  \uses{def:physics:phyx-mini-0471:phyxminiproblems-problemphyxmini0471-statelabel, def:physics:phyx-mini-0471:phyxminiproblems-problemphyxmini0471-processleg, def:physics:phyx-mini-0471:phyxminiproblems-problemphyxmini0471-segmentorientation, def:physics:phyx-mini-0471:phyxminiproblems-problemphyxmini0471-figureaxis, def:physics:phyx-mini-0471:phyxminiproblems-problemphyxmini0471-axisquantity, def:physics:phyx-mini-0471:phyxminiproblems-problemphyxmini0471-axisdisplayunit, def:physics:phyx-mini-0471:phyxminiproblems-problemphyxmini0471-axissymbol}
  Qualitative transcription of image 471.png. Physical state coordinates are stored separately in RectangularPVProcessSetup; this structure records the axes, units, labels, segment orientations, and directed arrows
\end{definition}
\begin{proof}
  This is the declared model definition of Pressure Volume Figure.
\end{proof}
\begin{definition}[Thermodynamic System Boundary]
  \label{def:physics:phyx-mini-0471:phyxminiproblems-problemphyxmini0471-thermodynamicsystemboundary}
  \lean{PhyXMiniProblems.ProblemPhyXMini0471.ThermodynamicSystemBoundary}
  Closed-system interpretation used for the fixed thermodynamic system
\end{definition}
\begin{proof}
  This is the declared model definition of Thermodynamic System Boundary.
\end{proof}
\begin{definition}[Rectangular PVProcess Setup]
  \label{def:physics:phyx-mini-0471:phyxminiproblems-problemphyxmini0471-rectangularpvprocesssetup}
  \lean{PhyXMiniProblems.ProblemPhyXMini0471.RectangularPVProcessSetup}
  \uses{def:physics:phyx-mini-0471:phyxminiproblems-problemphyxmini0471-volumequantity, def:physics:phyx-mini-0471:phyxminiproblems-problemphyxmini0471-pressurequantity, def:physics:phyx-mini-0471:phyxminiproblems-problemphyxmini0471-energyquantity, def:physics:phyx-mini-0471:phyxminiproblems-problemphyxmini0471-statelabel, def:physics:phyx-mini-0471:phyxminiproblems-problemphyxmini0471-processleg, def:physics:phyx-mini-0471:phyxminiproblems-problemphyxmini0471-processpath, def:physics:phyx-mini-0471:phyxminiproblems-problemphyxmini0471-processkind, def:physics:phyx-mini-0471:phyxminiproblems-problemphyxmini0471-pressurevolumefigure, def:physics:phyx-mini-0471:phyxminiproblems-problemphyxmini0471-thermodynamicsystemboundary}
  Independent physical observables of the system and its displayed processes. Internal energy is state-dependent, while heat and work are process-dependent. In particular, heat along acd is not defined from the recorded answer
\end{definition}
\begin{proof}
  This is the declared model definition of Rectangular PVProcess Setup.
\end{proof}
\begin{definition}[Matches Closed System Scenario]
  \label{def:physics:phyx-mini-0471:phyxminiproblems-problemphyxmini0471-matchesclosedsystemscenario}
  \lean{PhyXMiniProblems.ProblemPhyXMini0471.MatchesClosedSystemScenario}
  \uses{def:physics:phyx-mini-0471:phyxminiproblems-problemphyxmini0471-rectangularpvprocesssetup}
  Qualitative closed-system conditions implicit in the p--V exercise
\end{definition}
\begin{proof}
  This is the declared model definition of Matches Closed System Scenario.
\end{proof}
\begin{definition}[Matches Primary Pressure Volume Figure]
  \label{def:physics:phyx-mini-0471:phyxminiproblems-problemphyxmini0471-matchesprimarypressurevolumefigure}
  \lean{PhyXMiniProblems.ProblemPhyXMini0471.MatchesPrimaryPressureVolumeFigure}
  \uses{def:physics:phyx-mini-0471:phyxminiproblems-problemphyxmini0471-volumeincubicmetres, def:physics:phyx-mini-0471:phyxminiproblems-problemphyxmini0471-pressureinpascals, def:physics:phyx-mini-0471:phyxminiproblems-problemphyxmini0471-rectangularpvprocesssetup}
  Exact transcription of the primary p--V image. It records p\_a = p\_c = 3.0 × 10⁴ Pa and p\_b = p\_d = 8.0 × 10⁴ Pa; V\_a = V\_b = 2.0 × 10⁻³ m³ and V\_c = V\_d = 5.0 × 10⁻³ m³; vertical isochoric legs a → b and c → d; and horizontal isobaric legs a → c and b → d. No heat value, work value, or internal-energy value occurs in this structure
\end{definition}
\begin{proof}
  This is the declared model definition of Matches Primary Pressure Volume Figure.
\end{proof}
\begin{definition}[Matches Given Heat Readouts]
  \label{def:physics:phyx-mini-0471:phyxminiproblems-problemphyxmini0471-matchesgivenheatreadouts}
  \lean{PhyXMiniProblems.ProblemPhyXMini0471.MatchesGivenHeatReadouts}
  \uses{def:physics:phyx-mini-0471:phyxminiproblems-problemphyxmini0471-energyinjoules, def:physics:phyx-mini-0471:phyxminiproblems-problemphyxmini0471-rectangularpvprocesssetup}
  The two heat transfers stated in the prose. The heat along acd is absent: it is the current target rather than a datum
\end{definition}
\begin{proof}
  This is the declared model definition of Matches Given Heat Readouts.
\end{proof}
\begin{definition}[Has Physical Pressure Volume Coordinates]
  \label{def:physics:phyx-mini-0471:phyxminiproblems-problemphyxmini0471-hasphysicalpressurevolumecoordinates}
  \lean{PhyXMiniProblems.ProblemPhyXMini0471.HasPhysicalPressureVolumeCoordinates}
  \uses{def:physics:phyx-mini-0471:phyxminiproblems-problemphyxmini0471-volumeincubicmetres, def:physics:phyx-mini-0471:phyxminiproblems-problemphyxmini0471-pressureinpascals, def:physics:phyx-mini-0471:phyxminiproblems-problemphyxmini0471-rectangularpvprocesssetup}
  Positivity conditions selecting physical p--V coordinates
\end{definition}
\begin{proof}
  This is the declared model definition of Has Physical Pressure Volume Coordinates.
\end{proof}
\begin{definition}[Satisfies Quasistatic Boundary Work Law]
  \label{def:physics:phyx-mini-0471:phyxminiproblems-problemphyxmini0471-satisfiesquasistaticboundaryworklaw}
  \lean{PhyXMiniProblems.ProblemPhyXMini0471.SatisfiesQuasistaticBoundaryWorkLaw}
  \uses{def:physics:phyx-mini-0471:phyxminiproblems-problemphyxmini0471-volumeincubicmetres, def:physics:phyx-mini-0471:phyxminiproblems-problemphyxmini0471-pressureinpascals, def:physics:phyx-mini-0471:phyxminiproblems-problemphyxmini0471-energyinjoules, def:physics:phyx-mini-0471:phyxminiproblems-problemphyxmini0471-rectangularpvprocesssetup}
  Quasistatic pressure--volume boundary work. An isochoric leg does no work; on an isobaric leg the work done by the system is p (V\_final - V\_initial). This law is uniform over all legs and contains no requested heat value
\end{definition}
\begin{proof}
  This is the declared model definition of Satisfies Quasistatic Boundary Work Law.
\end{proof}
\begin{definition}[Satisfies Path Work Additivity]
  \label{def:physics:phyx-mini-0471:phyxminiproblems-problemphyxmini0471-satisfiespathworkadditivity}
  \lean{PhyXMiniProblems.ProblemPhyXMini0471.SatisfiesPathWorkAdditivity}
  \uses{def:physics:phyx-mini-0471:phyxminiproblems-problemphyxmini0471-energyinjoules, def:physics:phyx-mini-0471:phyxminiproblems-problemphyxmini0471-rectangularpvprocesssetup}
  Work on either complete route is the sum of the work on its two legs
\end{definition}
\begin{proof}
  This is the declared model definition of Satisfies Path Work Additivity.
\end{proof}
\begin{definition}[Satisfies Closed System First Law]
  \label{def:physics:phyx-mini-0471:phyxminiproblems-problemphyxmini0471-satisfiesclosedsystemfirstlaw}
  \lean{PhyXMiniProblems.ProblemPhyXMini0471.SatisfiesClosedSystemFirstLaw}
  \uses{def:physics:phyx-mini-0471:phyxminiproblems-problemphyxmini0471-energyinjoules, def:physics:phyx-mini-0471:phyxminiproblems-problemphyxmini0471-rectangularpvprocesssetup}
  The closed-system first law on each displayed leg and each complete path. Internal energy is a function of state, so the two complete paths share the same endpoint change. Neither field prescribes a numerical heat on acd
\end{definition}
\begin{proof}
  This is the declared model definition of Satisfies Closed System First Law.
\end{proof}
\begin{definition}[Answer Choice]
  \label{def:physics:phyx-mini-0471:phyxminiproblems-problemphyxmini0471-answerchoice}
  \lean{PhyXMiniProblems.ProblemPhyXMini0471.AnswerChoice}
  Labels of the four answers printed with the problem
\end{definition}
\begin{proof}
  This is the declared model definition of Answer Choice.
\end{proof}
\begin{definition}[displayed Heat In Joules]
  \label{def:physics:phyx-mini-0471:phyxminiproblems-problemphyxmini0471-answerchoice-displayedheatinjoules}
  \lean{PhyXMiniProblems.ProblemPhyXMini0471.AnswerChoice.displayedHeatInJoules}
  \uses{def:physics:phyx-mini-0471:phyxminiproblems-problemphyxmini0471-answerchoice}
  Heat in joules printed beside each answer label
\end{definition}
\begin{proof}
  This is the declared model definition of displayed Heat In Joules.
\end{proof}
\begin{definition}[recorded Dataset Answer]
  \label{def:physics:phyx-mini-0471:phyxminiproblems-problemphyxmini0471-recordeddatasetanswer}
  \lean{PhyXMiniProblems.ProblemPhyXMini0471.recordedDatasetAnswer}
  \uses{def:physics:phyx-mini-0471:phyxminiproblems-problemphyxmini0471-answerchoice}
  Dataset answer metadata, kept separate from all theorem premises
\end{definition}
\begin{proof}
  This is the declared model definition of recorded Dataset Answer.
\end{proof}
\begin{definition}[Matches Displayed Heat]
  \label{def:physics:phyx-mini-0471:phyxminiproblems-problemphyxmini0471-matchesdisplayedheat}
  \lean{PhyXMiniProblems.ProblemPhyXMini0471.MatchesDisplayedHeat}
  \uses{def:physics:phyx-mini-0471:phyxminiproblems-problemphyxmini0471-energyinjoules, def:physics:phyx-mini-0471:phyxminiproblems-problemphyxmini0471-rectangularpvprocesssetup, def:physics:phyx-mini-0471:phyxminiproblems-problemphyxmini0471-answerchoice}
  The modeled heat on acd agrees with the value printed for a choice
\end{definition}
\begin{proof}
  This is the declared model definition of Matches Displayed Heat.
\end{proof}
% --- Archon declaration topology end ---
% --- Archon physics formalization source end ---
